\documentclass[aspectratio=169,10pt]{beamer}
\usetheme{UU}

\usepackage[brazil]{babel}
\usepackage[utf8]{inputenc}
\usepackage{booktabs}
\usepackage{tabularx}
\usepackage{fontawesome5}
\usepackage{amsmath}
\usepackage{amssymb}

% --- Payoff ITA ---
\uupayoff{%
  {\rmfamily\mdseries\fontsize{11}{12}\selectfont Instituto Tecnológico de Aeronáutica}\\
  {\sffamily\fontseries{l}\small\itshape PO-233 Aprendizado de Máquina}%
}

% --- Metadados ---
\title[Recomendação de Filmes com SVD]{%
  Sistema de Recomendação de Filmes\\
  com Fatoração de Matrizes\\
  e Centralização por Usuários%
}
\author[D.\,Soares \and E.\,Costa]{%
  Denilson Gustavo de Araújo Soares\\[2pt]
  Enzo Mitsuo Tokushige Costa%
}
\institute{%
  Instituto Tecnológico de Aeronáutica\\[2pt]
  {\footnotesize Profª Ana Carolina Lorena}%
}
\date{Junho de 2026}

% ============================================================
\begin{document}
% ============================================================

% -------------------------------------------------------
% Slide 1 — Capa
% -------------------------------------------------------
\begin{frame}[plain]
  \titlepage
\end{frame}

% -------------------------------------------------------
% Slide 2 — Agenda
% -------------------------------------------------------
\begin{frame}{Agenda}
  \tableofcontents
\end{frame}

% ============================================================
\section{Motivação}
% ============================================================

\begin{frame}{Motivação}
  \begin{columns}[T]
    \begin{column}{0.52\textwidth}
      \begin{block}{O problema de recomendação}
        \small
        \begin{itemize}
          \item Plataformas como Netflix e Amazon dependem de recomendações relevantes
          \item Usuário avaliou, em média, \textbf{0,26\%} dos filmes disponíveis
          \item Esparsidade de \textbf{99,74\%} — similaridade direta é inviável
        \end{itemize}
      \end{block}
      \vskip 3pt
      \begin{block}{Abordagem adotada}
        \small\textbf{Fatoração de Matrizes} (SVD): vetores latentes que capturam preferências implícitas.
      \end{block}
    \end{column}
    \begin{column}{0.45\textwidth}
      \begin{alertblock}{O desafio central}
        \centering
        \vskip 4pt
        162.541 usuários\\
        $\times$\\
        59.047 filmes\\[6pt]
        $\approx$ 9,6 bilhões de pares\\[6pt]
        \textbf{apenas 25 milhões avaliados}
        \vskip 4pt
      \end{alertblock}
      \vskip 6pt
      \begin{exampleblock}{Meta}
        \small Estimar a nota que um usuário daria a filmes que \textbf{nunca assistiu} e recomendar os mais prováveis.
      \end{exampleblock}
    \end{column}
  \end{columns}
\end{frame}

% ============================================================
\section{Dataset — MovieLens 25M}
% ============================================================

\begin{frame}[shrink=5]{Dataset — MovieLens 25M}
  \begin{columns}[T]
    \begin{column}{0.50\textwidth}
      \begin{block}{Visão geral}
        \footnotesize
        \begin{tabular*}{\dimexpr\linewidth-12pt\relax}{@{\extracolsep{\fill}}lr}
          \toprule
          \textbf{Atributo} & \textbf{Valor} \\
          \midrule
          Avaliações        & 25.000.095 \\
          Usuários únicos   & 162.541    \\
          Filmes únicos     & 59.047     \\
          Esparsidade       & 99,74\%    \\
          \bottomrule
        \end{tabular*}
      \end{block}
      \vskip 3pt
      \begin{block}{Escala de notas}
        \footnotesize
        \begin{tabular*}{\dimexpr\linewidth-12pt\relax}{@{\extracolsep{\fill}}lr}
          \toprule
          \textbf{Propriedade} & \textbf{Valor} \\
          \midrule
          Intervalo        & 0,5 a 5,0 (passo 0,5) \\
          Média / Mediana  & 3,53 \;/\; 3,50       \\
          Notas $\geq$ 3,5 & 62,5\%                \\
          \bottomrule
        \end{tabular*}
      \end{block}
    \end{column}
    \begin{column}{0.47\textwidth}
      \begin{block}{Colunas do arquivo \texttt{ratings.csv}}
        \footnotesize
        \begin{tabular*}{\dimexpr\linewidth-12pt\relax}{@{\extracolsep{\fill}}ll}
          \toprule
          \textbf{Coluna} & \textbf{Descrição} \\
          \midrule
          \texttt{userId}    & Identificador do usuário \\
          \texttt{movieId}   & Identificador do filme   \\
          \texttt{rating}    & Nota dada (0,5 – 5,0)   \\
          \texttt{timestamp} & Data da avaliação (Unix) \\
          \bottomrule
        \end{tabular*}
      \end{block}
      \vskip 3pt
      \begin{block}{Colunas do arquivo \texttt{movies.csv}}
        \footnotesize
        \begin{tabular*}{\dimexpr\linewidth-12pt\relax}{@{\extracolsep{\fill}}ll}
          \toprule
          \textbf{Coluna} & \textbf{Descrição} \\
          \midrule
          \texttt{movieId} & Identificador do filme          \\
          \texttt{title}   & Título e ano                    \\
          \texttt{genres}  & Gêneros separados por \texttt{|} \\
          \bottomrule
        \end{tabular*}
      \end{block}
    \end{column}
  \end{columns}
\end{frame}

% ============================================================
\section{Pré-processamento}
% ============================================================

\begin{frame}[shrink=3]{Pré-processamento}
  \begin{columns}[T]
    \begin{column}{0.52\textwidth}
      \begin{block}{Pipeline de pré-processamento}
        \small
        \begin{enumerate}\itemsep2pt
          \item Remoção de duplicatas \textbf{(0 encontradas)}
          \item Conversão \texttt{timestamp} $\rightarrow$ \texttt{datetime}
          \item Remoção de filmes sem avaliações \textbf{(3.376 removidos)}
          \item Filtro: usuários $\geq$ 20 avaliações; filmes $\geq$ 50
        \end{enumerate}
      \end{block}
      \vskip 3pt
      \begin{block}{Divisão treino / teste — \emph{split} temporal}
        \small
        \begin{tabular*}{\dimexpr\linewidth-12pt\relax}{@{\extracolsep{\fill}}lrr}
          \toprule
          \textbf{Conjunto} & \textbf{Ratings} & \textbf{Até} \\
          \midrule
          Treino (80\%) & 19.715.942 & 2016-05-28 \\
          Teste  (20\%) &  4.928.986 & depois     \\
          \bottomrule
        \end{tabular*}
      \end{block}
    \end{column}
    \begin{column}{0.45\textwidth}
      \begin{alertblock}{Por que \emph{split} temporal?}
        \small Simula o cenário real: o modelo só vê avaliações \textbf{passadas} para recomendar filmes futuros. Evita vazamento de informação.
      \end{alertblock}
      \vskip 3pt
      \begin{block}{Após filtragem}
        \small
        \begin{tabular*}{\dimexpr\linewidth-12pt\relax}{@{\extracolsep{\fill}}lr}
          \toprule
          \textbf{Atributo} & \textbf{Valor} \\
          \midrule
          Usuários & 162.540    \\
          Filmes   & 13.176     \\
          Ratings  & 24.644.928 \\
          \bottomrule
        \end{tabular*}
      \end{block}
    \end{column}
  \end{columns}
\end{frame}

% ============================================================
\section{Fundamentação: SVD e Centralização}
% ============================================================

\begin{frame}{Fatoração de Matrizes com SVD}
  \begin{columns}[T]
    \begin{column}{0.52\textwidth}
      \begin{block}{Ideia central}
        \small A matriz esparsa $\mathbf{R}$ (usuários $\times$ filmes) é aproximada:
        \[ \mathbf{R} \approx \mathbf{U}\,\boldsymbol{\Sigma}\,\mathbf{V}^\top \]
        \vskip -4pt
        \begin{itemize}
          \item $\mathbf{U}$: perfil latente de cada usuário
          \item $\mathbf{V}$: perfil latente de cada filme
          \item $\boldsymbol{\Sigma}$: importância de cada fator
        \end{itemize}
      \end{block}
      \vskip 3pt
      \begin{block}{Predição de nota}
        \small
        \[ \hat{r}_{ui} = \mu_u + \mathbf{u}_u^\top\,\mathbf{v}_i \]
        \vskip -4pt
        onde $\mu_u$ é a média do usuário $u$.
      \end{block}
    \end{column}
    \begin{column}{0.45\textwidth}
      \begin{block}{Parâmetros do modelo}
        \small
        \begin{tabular*}{\dimexpr\linewidth-12pt\relax}{@{\extracolsep{\fill}}lr}
          \toprule
          \textbf{Parâmetro} & \textbf{Valor} \\
          \midrule
          Fatores latentes ($k$) & 50      \\
          Variância explicada    & 18,10\% \\
          Esparsidade da matriz  & 98,83\% \\
          \bottomrule
        \end{tabular*}
      \end{block}
      \vskip 3pt
      \begin{alertblock}{Centralização por usuário}
        \small Antes do SVD, subtrai-se $\mu_u$ de cada nota de $u$:
        \[ \tilde{r}_{ui} = r_{ui} - \mu_u \]
        \vskip -4pt
        Elimina o viés de usuários que sistematicamente dão notas altas ou baixas.
      \end{alertblock}
    \end{column}
  \end{columns}
\end{frame}

% ============================================================
\section{Método de Recomendação}
% ============================================================

\begin{frame}[shrink=8]{Similaridade de Cosseno entre Itens}
  \begin{columns}[T]
    \begin{column}{0.52\textwidth}
      \begin{block}{Vetores latentes dos filmes}
        \small SVD produz \texttt{item\_factors}: matriz $(13.176 \times 50)$.\\[3pt]
        Similaridade entre filmes $i$ e $j$:
        \[
          \text{sim}(i,j) = \frac{\mathbf{v}_i^\top\mathbf{v}_j}{\|\mathbf{v}_i\|\|\mathbf{v}_j\|} \in [0,1]
        \]
      \end{block}
      \vskip 3pt
      \begin{block}{Geração de recomendações}
        \small
        \begin{enumerate}\itemsep2pt
          \item Busca filmes avaliados pelo usuário no treino
          \item Para cada filme, encontra os 20 mais similares
          \item Pondera similaridade pela nota do usuário
          \item Remove já vistos e retorna Top-10
        \end{enumerate}
      \end{block}
    \end{column}
    \begin{column}{0.45\textwidth}
      \begin{block}{Exemplo de similaridade}
        \footnotesize Filme base: \emph{A Fish Called Wanda} (1988)
        \vskip 2pt
        \begin{tabular*}{\dimexpr\linewidth-12pt\relax}{@{\extracolsep{\fill}}lr}
          \toprule
          \textbf{Filme similar} & \textbf{sim} \\
          \midrule
          Monty Python: And Now\ldots  & 0,796 \\
          Monty Python's Life of Brian & 0,764 \\
          Monty Python: Hollywood Bowl & 0,735 \\
          Monty Python: Holy Grail     & 0,710 \\
          The Commitments (1991)       & 0,702 \\
          \bottomrule
        \end{tabular*}
        \vskip 4pt
        \footnotesize\textit{Filmes do mesmo universo emergem \textbf{sem supervisão explícita} de gênero.}
      \end{block}
    \end{column}
  \end{columns}
\end{frame}

% ============================================================
\section{Resultados}
% ============================================================

\begin{frame}{Resultados — RMSE e MAE}
  \begin{columns}[T]
    \begin{column}{0.52\textwidth}
      \begin{block}{Predição de nota}
        \small
        \begin{tabular*}{\dimexpr\linewidth-12pt\relax}{@{\extracolsep{\fill}}lcc}
          \toprule
          \textbf{Métrica} & \textbf{Sem centr.} & \textbf{Com centr.} \\
          \midrule
          RMSE & 3,4339 & \textbf{0,9635} \\
          MAE  & 3,2834 & \textbf{0,7177} \\
          \bottomrule
        \end{tabular*}
      \end{block}
      \vskip 3pt
      \begin{block}{Baseline ingênuo (média global)}
        \small $\text{RMSE}_{\text{baseline}} \approx 1{,}0221$ — predizer sempre a média global.
      \end{block}
    \end{column}
    \begin{column}{0.45\textwidth}
      \begin{alertblock}{Impacto da centralização}
        \small A centralização reduziu o RMSE de \textbf{3,43 para 0,96} — queda de \textbf{72\%}.\\[2pt]
        Sem ela, o modelo ignorava o viés individual de cada usuário.
      \end{alertblock}
      \vskip 3pt
      \begin{exampleblock}{Comparação com baseline}
        \small O modelo supera o baseline ingênuo:\\[2pt]
        $0{,}9635 < 1{,}0221$ \quad (\textbf{RMSE})\\
        $0{,}7177 < 1{,}06$ \quad (\textbf{MAE})
      \end{exampleblock}
    \end{column}
  \end{columns}
\end{frame}

\begin{frame}{Resultados — Precision@10}
  \begin{columns}[T]
    \begin{column}{0.52\textwidth}
      \begin{block}{Definição da métrica}
        \small
        \[ \text{Precision@}k = \frac{|\text{top-}k \cap \text{relevantes}|}{k} \]
        \vskip -4pt
        \begin{itemize}
          \item \textbf{Candidatos:} filmes avaliados no teste e conhecidos pelo modelo
          \item \textbf{Relevante:} nota $\geq$ 4,0 \quad ($k = 10$, 200 usuários)
        \end{itemize}
      \end{block}
      \vskip 3pt
      \begin{block}{Resultado}
        \small\centering
        $\text{Precision@10} = \mathbf{0{,}5743}$\\[2pt]
        Entre os 10 sugeridos, \textbf{57,4\%} são relevantes.
      \end{block}
    \end{column}
    \begin{column}{0.45\textwidth}
      \begin{alertblock}{Interpretação do protocolo}
        \small Os candidatos são restritos aos filmes que o usuário \textbf{de fato avaliou} no teste — não o catálogo inteiro.\\[2pt]
        Torna a métrica equivalente a um \emph{ranking} entre opções conhecidas.
      \end{alertblock}
      \vskip 3pt
      \begin{exampleblock}{Leitura prática}
        \small Em cada lista de 10, o modelo acerta \textbf{5 a 6 filmes} que o usuário consideraria bons — sem ter visto essas avaliações no treino.
      \end{exampleblock}
    \end{column}
  \end{columns}
\end{frame}

\begin{frame}{Exemplo de Recomendação}
  \begin{columns}[T]
    \begin{column}{0.48\textwidth}
      \begin{block}{Histórico do usuário 1 (nota 5,0)}
        \small
        \begin{tabular}{l}
          \toprule
          \textbf{Filmes avaliados} \\
          \midrule
          Lost in Translation (2003) \\
          Requiem for a Dream (2000) \\
          Three Colors: Blue (1993) \\
          The Seventh Seal (1957) \\
          Look at Me (2004) \\
          \bottomrule
        \end{tabular}
      \end{block}
    \end{column}
    \begin{column}{0.50\textwidth}
      \begin{block}{Top-5 recomendações geradas}
        \footnotesize
        \begin{tabular*}{\dimexpr\linewidth-12pt\relax}{@{\extracolsep{\fill}}ll}
          \toprule
          \textbf{Filme recomendado} & \textbf{Gênero} \\
          \midrule
          Discreet Charm\ldots (1972) & Comedy/Drama \\
          Vivre sa vie (1962)         & Drama \\
          Andrei Rublev (1969)        & Drama/War \\
          Viridiana (1961)            & Comedy/Drama \\
          Last Year at Marienbad      & Drama/Mystery \\
          \bottomrule
        \end{tabular*}
      \end{block}
      \vskip 4pt
      \begin{exampleblock}{Interpretação}
        \small O modelo inferiu preferência por \textbf{cinema de arte europeu} sem nenhuma informação explícita de gênero — apenas pelos padrões latentes da matriz.
      \end{exampleblock}
    \end{column}
  \end{columns}
\end{frame}

% ============================================================
\section{Limitações e Trabalhos Futuros}
% ============================================================

\begin{frame}{Limitações e Trabalhos Futuros}
  \begin{columns}[T]
    \begin{column}{0.50\textwidth}
      \begin{alertblock}{Limitações identificadas}
        \begin{itemize}
          \item \textbf{Cold start:} usuários e filmes novos não têm vetor latente
          \item \textbf{SVD estático:} não atualiza incrementalmente com novas avaliações
          \item \textbf{18,10\% de variância explicada} — 50 componentes podem ser insuficientes
          \item Precision@10 avaliada apenas em filmes \emph{já vistos} no teste, não no catálogo completo
        \end{itemize}
      \end{alertblock}
    \end{column}
    \begin{column}{0.47\textwidth}
      \begin{block}{Possíveis extensões}
        \begin{enumerate}
          \item Substituir SVD por \textbf{ALS} (\emph{Alternating Least Squares}) — mais adequado para dados implícitos
          \item Incorporar \textbf{conteúdo} (gêneros, ano, tags) para mitigar cold start
          \item Avaliar com \emph{leave-one-out} sobre o catálogo completo
          \item Aumentar $k$ e comparar variância explicada
        \end{enumerate}
      \end{block}
    \end{column}
  \end{columns}
\end{frame}

% -------------------------------------------------------
% Encerramento
% -------------------------------------------------------
\thankframe{Obrigado! \quad Perguntas?}{%
  Denilson Gustavo de Araújo Soares\\
  Enzo Mitsuo Tokushige Costa\\[4pt]
  \texttt{PO-233 Aprendizado de Máquina}%
}

\end{document}
